% ===========================================================================
% SVNN Theoretical Completeness: Theorem 5 (Computational Necessity) + Theorem 6 (Generalization)
% ===========================================================================
% Insert after % ===========================================================================
% III.X. Chebyshev-KAN Satisfies the SVNN Conditions
% ===========================================================================
% To be inserted into section_svnn.tex after \subsection{Standard MLPs Do Not Admit Tight SVNN Error Bounds}
% and before \subsection{Implications for Compiler Design}

\subsection{Chebyshev-KAN Satisfies the SVNN Conditions}
\label{sec:svnn-chebykan}

Having established that B-spline KAN satisfies the SVNN conditions
({\S}\ref{sec:svnn-kan}) and that standard MLPs do not
({\S}\ref{sec:svnn-mlp}), we now demonstrate that the SVNN framework
extends to alternative basis function families. We prove that KAN
architectures using \textit{Chebyshev polynomial} basis functions---a
globally-supported, orthogonal polynomial basis---also satisfy Conditions~1
and~2, establishing that SVNN membership is not tied to the local-support
property of B-splines but rather to the structural decomposition and
curvature computability that both bases provide.

\vspace{4pt}
\begin{proposition}[Chebyshev-KAN Is Structurally Verifiable]
\label{prop:chebykan-svnn}
A KAN architecture using Chebyshev polynomial basis functions
(ChebyKAN)~\cite{bozorgasl2024chebykan}, defined as
\begin{equation}
\phi_{j,i}(x) = \sum_{n=0}^{N} c_{j,i,n} \cdot T_n(\tanh(x))
\label{eq:chebykan_def}
\end{equation}
where $T_n$ is the Chebyshev polynomial of degree $n$ and $c_{j,i,n}$ are
learnable coefficients, satisfies SVNN Conditions~1 and~2 of Theorem~2.
Furthermore, ChebyKAN enjoys Z3 verifiability via polynomial NRA: the
Chebyshev basis functions $T_n$ are degree-$n$ polynomials expressible
directly in Z3's nonlinear real arithmetic theory, enabling compositional
verification without the segment enumeration required for B-spline basis.
\end{proposition}

\vspace{4pt}
\noindent\textit{Proof.} We verify each SVNN condition separately, then
establish Z3 verifiability.

\vspace{2pt}
\noindent\textit{Condition~1 (Operation-Type Closure).}
A ChebyKAN layer with $d_{\text{in}}$ inputs and $d_{\text{out}}$ outputs
decomposes as:
\begin{equation}
y_j = \sum_{i=1}^{d_{\text{in}}} \left[ \sum_{n=0}^{N} c_{j,i,n} \cdot
T_n(\tanh(x_i)) \right]
\label{eq:chebykan_layer}
\end{equation}

This computation naturally separates into three pure operation types:
\begin{enumerate}
    \item \textbf{Bound mapping (element-wise):}
          $u_i = \tanh(x_i)$ for each input $i$. The hyperbolic tangent
          maps $\mathbb{R} \to (-1,1)$, placing inputs in the natural
          domain of Chebyshev polynomials. This is an element-wise
          univariate operation (TanhLUT node in the IR).

    \item \textbf{Polynomial basis evaluation (element-wise):}
          For each pair $(i,n)$, compute $v_{i,n} = T_n(u_i)$. Since
          Chebyshev polynomials satisfy the recurrence
          $T_{n+1}(x) = 2x \cdot T_n(x) - T_{n-1}(x)$ with $T_0(x)=1$
          and $T_1(x)=x$, each $T_n$ is a univariate function of $u_i$.
          These are element-wise univariate operations (ChebyLUT nodes).

    \item \textbf{Linear combination:}
          $y_j = \sum_{i,n} c_{j,i,n} \cdot v_{i,n}$ aggregates the
          basis-evaluated values via learnable coefficients. This is a
          linear transformation (MatMul node with coefficient matrix $C$).
\end{enumerate}

Crucially, no operation mixes linear and nonlinear computations: the
nonlinearity is confined to steps~1--2 (element-wise tanh and polynomial
evaluation), while step~3 is purely linear. This decomposition satisfies
Condition~1. The IR representation is:
\begin{equation}
\text{Input} \xrightarrow{\text{TanhLUT}} u \xrightarrow{\text{ChebyLUT}}
v \xrightarrow{\text{MatMul}(C)} \text{Output}
\end{equation}
where each node is provably of a single pure type (element-wise or linear).

\vspace{2pt}
\noindent\textit{Condition~2 (Univariate Curvature Bound).}
We must show that the second derivative $f''(x)$ of each composite
activation $f(x) = \sum_n c_n \cdot T_n(\tanh(x))$ is bounded and
computable from the architecture parameters $(N, \{c_n\})$ alone.

\textbf{Step 2a (Tanh bound mapping):}
The function $h(x) = \tanh(x)$ is $C^\infty$ with bounded second
derivative. The analytical bound is:
\begin{equation}
\sup_{x \in \mathbb{R}} |\tanh''(x)| = \frac{2}{\sqrt{27}} \approx 0.385
\label{eq:tanh_m2}
\end{equation}
This is a universal constant, independent of any learnable parameters.

\textbf{Step 2b (Chebyshev polynomial curvature):}
Let $g(y) = \sum_{n=0}^{N} c_n \cdot T_n(y)$ be the polynomial of degree
$N$ evaluated on $y \in (-1,1)$. Chebyshev polynomials $T_n$ are
degree-$n$ polynomials satisfying $|T_n(y)| \leq 1$ for all $y \in
[-1,1]$~\cite{mason2002chebyshev}. By the Markov Brothers'
Inequality~\cite{markov1916limit}, for any polynomial $p$ of degree $N$
on $[-1,1]$:
\begin{equation}
\max_{y \in [-1,1]} |p'(y)| \leq N^2 \cdot \max_{y \in [-1,1]} |p(y)|
\label{eq:markov_first}
\end{equation}
\begin{equation}
\max_{y \in [-1,1]} |p''(y)| \leq \frac{N^2(N^2-1)}{3} \cdot
\max_{y \in [-1,1]} |p(y)|
\label{eq:markov_second}
\end{equation}

For the polynomial $g(y) = \sum_n c_n T_n(y)$, we have
$\max_{y \in [-1,1]} |g(y)| \leq \sum_n |c_n|$ (triangle inequality
with $|T_n(y)| \leq 1$). Applying Eq.~\ref{eq:markov_second}:
\begin{equation}
M_2^{\text{Cheby}}(g) := \max_{y \in [-1,1]} |g''(y)| \leq
\frac{N^2(N^2-1)}{3} \cdot \sum_{n=0}^{N} |c_n|
\label{eq:cheby_m2_bound}
\end{equation}

This bound is \textit{computable from parameters alone}: given the
polynomial degree $N$ and the learned coefficients $\{c_n\}$, the compiler
can compute the right-hand side of Eq.~\ref{eq:cheby_m2_bound} without
executing the model on any inputs. This satisfies the ``parameter-computable
curvature'' requirement of Condition~2.

\textbf{Step 2c (Composition via chain rule):}
The full activation is $f(x) = g(\tanh(x))$. Applying the chain rule:
\begin{equation}
f''(x) = g''(\tanh x) \cdot \text{sech}^4(x) - 2g'(\tanh x) \cdot
\tanh(x) \cdot \text{sech}^2(x)
\label{eq:cheby_f_second}
\end{equation}

Since $|\text{sech}(x)| \leq 1$ and $|\tanh(x)| < 1$ for all $x$, we
bound:
\begin{equation}
|f''(x)| \leq |g''(\tanh x)| + 2|g'(\tanh x)| \leq M_2^{\text{Cheby}}(g)
+ 2 M_1^{\text{Cheby}}(g)
\label{eq:cheby_f_m2_total}
\end{equation}
where $M_1^{\text{Cheby}}(g) \leq N^2 \cdot \sum_n |c_n|$ by
Eq.~\ref{eq:markov_first}. Substituting Eq.~\ref{eq:cheby_m2_bound}:
\begin{equation}
M_2(f) \leq \left[\frac{N^2(N^2-1)}{3} + 2N^2\right] \cdot \sum_{n=0}^{N}
|c_n| = \frac{N^2(N^2+5)}{3} \cdot \sum_{n=0}^{N} |c_n|
\label{eq:cheby_final_m2}
\end{equation}

Equation~\ref{eq:cheby_final_m2} provides an \textit{a priori} curvature
bound computable from the polynomial degree $N$ (an architectural
hyperparameter) and the learned coefficients $\{c_{j,i,n}\}$ (model
parameters). The compiler can evaluate this bound for every activation in
the network before generating SCL code, satisfying Condition~2.

\vspace{2pt}
\noindent\textit{Z3 Verifiability.}
The compositional verification pipeline ({\S}\ref{sec:compositional}) requires that each
activation function can be abstracted as a finite set of polynomial
constraints verifiable by Z3's SMT solver. ChebyKAN satisfies this
requirement in two ways:

\begin{enumerate}
    \item \textbf{Polynomial basis is NRA-native.} Chebyshev polynomials
          $T_n(y)$ are expressible as explicit polynomials in $y$ (e.g.,
          $T_2(y) = 2y^2 - 1$, $T_3(y) = 4y^3 - 3y$). Z3's Nonlinear
          Real Arithmetic (NRA) theory handles polynomial equalities and
          inequalities natively~\cite{de2008z3}, making each segment's
          verification a polynomial feasibility query.

    \item \textbf{Tanh preprocessing is standard.} The tanh bound mapping
          $u = \tanh(x)$ is treated identically to B-spline's SiLU base
          path: a $C^\infty$ univariate function approximated via LUT,
          whose error bound is $M_{\tanh} h^2/8$ (Eq.~\ref{eq:tanh_m2}).
          The LUT approximation introduces a bounded error that propagates
          through the polynomial evaluation, and the combined constraint
          system remains polynomial in the LUT-approximated $u$ variable.
\end{enumerate}

In contrast to B-spline basis functions---which require segment-by-segment
enumeration due to their local support---Chebyshev polynomials are
\textit{globally supported}: each $T_n$ is defined on the entire interval
$(-1,1)$. This eliminates the segment enumeration overhead but requires
computing bounds over the full domain. The tradeoff is: B-spline achieves
tighter per-segment bounds via the segment-aware $M_2$ refinement
({\S}\ref{sec:segment-aware}), while ChebyKAN uses global polynomial
bounds (Eq.~\ref{eq:cheby_m2_bound}) that are looser but analytically
simpler. Both approaches yield finite, computable bounds---the SVNN
requirement.

Empirical Z3 verification confirms the theoretical prediction: a
ChebyKAN$[28,16,4]$ with polynomial degree $N=5$ achieves 496/512
Z3-verifiable components (96.9\%, E54, {\S}\ref{sec:cross_domain}), demonstrating that the
polynomial NRA encoding successfully produces verifiable constraints for
the majority of activations. The Markov-based analytical $M_2$ bound
(Eq.~\ref{eq:cheby_m2_bound}) is $5.4\times$ looser than the empirical
$M_2$ (mean $8.5$ vs.\ $45.6$ in Layer~0, E54), confirming that the
worst-case nature of Markov's inequality is the primary source of
non-verifiability. The 16 unverifiable components (3.1\%) correspond to
activations where this global bound exceeds Z3's numeric tolerance---a
fundamental limitation of any global polynomial bound rather than a
deficiency of ChebyKAN specifically.

\vspace{2pt}
\noindent\textbf{Contrast with MLP+Tanh (Structural vs.\ Activation-Based Distinction).}
Table~\ref{tab:svnn_taxonomy} shows that MLP+Tanh lies on the
\textit{exterior} (violates Condition~1) despite using the same tanh
activation that ChebyKAN uses for bound mapping. This contrast isolates
the structural nature of the SVNN conditions: it is not the activation
function itself that determines SVNN membership, but rather how the
architecture \textit{decomposes} that activation relative to linear
transformations.

\begin{itemize}
    \item \textbf{MLP+Tanh:} The standard MLP layer computes
          $\mathbf{h} = \tanh(W\mathbf{x} + \mathbf{b})$, coupling the
          affine transformation $(W\mathbf{x} + \mathbf{b})$ and the
          nonlinear activation $\tanh(\cdot)$ in a single operation. There
          is no intermediate variable whose error can be separately
          bounded---the structural induction of Theorem~2 cannot be
          initiated. MLP+Tanh thus violates Condition~1.

    \item \textbf{ChebyKAN:} The ChebyKAN layer first applies tanh
          \textit{element-wise} to each input (step~1), then evaluates
          polynomial basis functions element-wise (step~2), and finally
          performs a linear combination (step~3). Each step is provably
          pure (either element-wise or linear), enabling the per-operation
          error induction. ChebyKAN satisfies Condition~1.
\end{itemize}

This distinction underscores the central insight of the SVNN framework:
\textit{architectures must be designed for verifiability}. The same
activation function (tanh) yields verifiable or non-verifiable behavior
depending on whether the architecture exposes the structural decomposition
required by Condition~1. The compiler cannot ``fix'' an architecture that
couples operations; it can only exploit the structure that the architecture
provides.

\vspace{2pt}
\noindent\textbf{Comparison with B-spline KAN (Local vs.\ Global Basis).}
Proposition~\ref{prop:kan-svnn} established that B-spline KAN satisfies the
SVNN conditions via \textit{segment-aware} curvature bounds: each B-spline
is piecewise-polynomial with $C^2$ continuity, and the compiler computes
per-segment $M_2^{(k)}$ values tightening the global bound by $6.0\times$
on average. Proposition~\ref{prop:chebykan-svnn} shows that ChebyKAN
satisfies the same conditions via \textit{Markov-based global bounds}: the
Chebyshev polynomial $g(y)$ has a single, analytically-computable $M_2$
bound (Eq.~\ref{eq:cheby_m2_bound}) that applies to the entire domain
$(-1,1)$.

The two bases represent complementary points on the tightness-complexity
spectrum:
\begin{itemize}
    \item \textbf{B-spline (local support):} Tighter bounds via
          segment-aware refinement, but requires $O(G)$ segment enumeration
          where $G$ is the grid size. The segment-aware $M_2^{(k)}$ bounds
          exploit the locality of B-spline support: outside a basis
          function's $k+1$ non-zero knot intervals, it contributes zero
          error.

    \item \textbf{Chebyshev (global support):} Single global bound via
          Markov's inequality (Eq.~\ref{eq:cheby_m2_bound}), computed in
          $O(1)$ for the entire activation. The bound is looser
          (worst-case over the full domain) but analytically simpler and
          requires no segment enumeration.
\end{itemize}

Both approaches yield \textit{computable} bounds from parameters alone,
satisfying Condition~2. The SVNN framework is thus robust to the basis
function choice: it accommodates both local-support (B-spline) and
global-support (Chebyshev) polynomial bases, as long as the curvature $M_2$
is analytically bounded and the structural decomposition (Condition~1)
holds. This generality confirms that SVNN conditions characterize a
fundamental architectural property, not an artifact specific to one basis
family. $\square$

\vspace{4pt}
\noindent\textbf{Remark 4 (Wavelet-KAN and Fourier-KAN).}
The proof technique of Proposition~\ref{prop:chebykan-svnn}---verifying
Condition~1 via explicit decomposition and Condition~2 via analytical
derivative bounds---extends naturally to other orthogonal polynomial and
wavelet bases. Wavelet-KAN~\cite{bozorgasl2024chebykan}, which uses
Daubechies or Haar wavelet basis functions, satisfies Condition~1 by the
same decomposition argument (element-wise wavelet evaluation + linear
combination). Condition~2 holds if the wavelet basis is
\textit{compactly-supported and piecewise-polynomial}---properties
satisfied by Daubechies wavelets but not by all wavelet families (e.g.,
Morlet wavelets involve Gaussian terms and violate Condition~2).
Fourier-KAN, using $\sin$ and $\cos$ basis functions, satisfies Condition~1
but faces the same Z3-undecidability issue as SiLU (transcendental
functions). Table~\ref{tab:svnn_taxonomy} marks Wavelet-KAN with
``$\sim$'' (depends on wavelet choice) to reflect this nuance: the SVNN
conditions are \textit{basis-specific}, not universally applicable to all
function approximation schemes.
 and before \subsection{Implications for Compiler Design}

\subsection{Verification-Cost Separation in the $\varepsilon$-Dimension}
\label{sec:svnn-necessity}

Theorem~\ref{thm:svnn_conditions} establishes that Conditions~1--2 are
\textit{sufficient} for structural verifiability. We now establish a
\emph{correct} complexity separation. The earlier claim---that tight-bound
computation is NP-hard for non-SVNN architectures while SVNN admits
polynomial-time bounds---compares an exact quantity (NP-hard for \emph{both}
classes) with an upper bound (polynomial for \emph{both} classes: interval
propagation and CROWN are polynomial-time for MLPs), so it does not
separate the classes. The valid separation operates in the
\emph{$\varepsilon$-dimension}: the \emph{decision problem itself} is
polynomial for SVNN and NP-hard for general coupled architectures, and
SVNN achieves the Pareto point (polynomial time, zero slack)
simultaneously---which no non-SVNN method can.

\setcounter{theorem}{4}
\begin{theorem}[$\varepsilon$-Separation of LUT-Verification Complexity]
\label{thm:necessity}
Let $\varepsilon$-Verify$(N, \mathrm{LUT}, \varepsilon)$ be the decision
problem: ``is $\sup_{x \in \Omega} |N(x) - \mathrm{LUT}(x)| \leq \varepsilon$?''
for a feedforward network $N$ on a compact box $\Omega$.

\begin{enumerate}
    \item \textbf{SVNN regime (polynomial).} For SVNNs with order-$k$
    LUTs, the sharp bound of Theorem~\ref{thm:sharp-lut} is computable in
    $O(L \cdot d_{\max}^2 \cdot N)$ design time
    ($N = \Theta(\varepsilon^{-1/k})$ cells to certify error $\leq
    \varepsilon$) and $O(L \cdot d_{\max}^2)$ evaluation time; the
    $\varepsilon$-decision is $O(1)$ given the LUT (compare the closed-form
    bound to $\varepsilon$). For B-spline KAN activations the order-$k$
    LUT reproduces the spline \emph{exactly} (polynomial reproduction:
    each segment is a cubic polynomial, degree-3 interpolation through 4
    points is exact---verified to $5 \times 10^{-8}$, E65), so
    $\varepsilon = 0$ is achieved at training-defined cost.

    \item \textbf{Non-SVNN regime (NP-hard, self-contained).} Let
    $N_\varphi$ be the network with product gates computing
    \begin{equation}
    F(x) \;=\; \sum_{j} \prod_{i \in C_j} \bigl(1 - \ell_{ij}(x)\bigr)
    \label{eq:product_gate}
    \end{equation}
    for a 3-SAT formula $\varphi$ with clauses $C_j$ and literals
    $\ell_{ij}$ affine in $x \in [0,1]^n$ (arity-3 coupled gates, one
    hidden layer, $O(n + m)$ neurons). Then $F(x) = 0$ if and only if
    $\varphi$ is satisfiable: $F(x^*) = 0$ forces every clause to have a
    literal with value exactly $1$; the induced assignment is consistent
    (no complementary pair can both be $1$---$\ell = x_v$ and
    $\ell' = 1 - x_v$ cannot both equal $1$), so $\varphi$ is satisfied;
    conversely a satisfying boolean assignment gives $F = 0$. Hence the
    exact problem $\varepsilon$-Verify with $\varepsilon = 0$ is
    NP-hard (3-SAT reduces in polynomial time with $O(n+m)$ neurons).

    \item \textbf{Gap version (NP-hard for every constant $\varepsilon_0$).}
    The published gap-preserving reductions for ReLU verification
    (Katz et al.~\cite{katz2017reluplex}; Tjeng et al.; S\"{a}lzer--Lange)
    produce networks with a constant separation between SAT and UNSAT
    cases, implying that $\varepsilon_0$-Verify is NP-hard for every fixed
    $\varepsilon_0 < 1$ (after normalization): deciding whether a
    poly-time-checkable piecewise-linear approximant is within any constant
    $\varepsilon_0$ of a general network is NP-hard. Under the
    Exponential-Time Hypothesis, exact verification of general networks
    requires $2^{\Omega(n)}$ time.

    \item \textbf{Pareto corollary.} For SVNNs the point
    $(\mathrm{time} = \mathrm{poly}, \mathrm{slack} = 0)$ is
    \emph{simultaneously achievable}; for general networks, any
    polynomial-time sound bound has additive slack $\geq \varepsilon_0$
    (else it would decide $\varepsilon_0$-Verify on the gap instances),
    and any slack-$0$ bound costs $2^{\Omega(n)}$ (ETH). No point of the
    non-SVNN frontier dominates any point of the SVNN frontier in both
    coordinates.
\end{enumerate}
\end{theorem}

\vspace{2pt}
\noindent\textit{Proof.}
\textbf{(1)} The closed-form bound of Theorem~\ref{thm:sharp-lut}
(Eq.~\ref{eq:sharp_soundness}) requires scanning $O(N)$ cells per
activation; there are $O(L \cdot d_{\max}^2)$ activations, giving
$O(L \cdot d_{\max}^2 \cdot N)$ design time; the evaluation is the
LUT lookups themselves, $O(L \cdot d_{\max}^2)$. Spline reproduction:
a cubic B-spline segment is a cubic polynomial; the order-4 LUT
(degree-3 interpolation through 4 points per window) reproduces it
exactly (verified numerically, E65). \textbf{(2)} The product-gate
encoding $F$: each clause contributes a non-negative term
$\prod_i (1 - \ell_{ij}(x)) \geq 0$, so $F(x) = 0$ iff every term
vanishes iff every clause has some literal $\ell_{ij}(x) = 1$. For
fractional $x^*$, the set of literals with value $1$ induces a
consistent boolean assignment: if $x^*_v = 1$ is required by one clause
and $x^*_v = 0$ by another, then $x^*_v = 1$ and $1 - x^*_v = 1$
simultaneously---impossible. Hence $\varphi$ is satisfiable; the
converse is immediate. The reduction is self-contained (verified on
200 random formulas with $n = 5$ variables, both directions, including
fractional-grid search, E65). \textbf{(3)} Standard gap-preserving
reductions (cited); ETH transfer is standard. \textbf{(4)} By
contrapositive: a polynomial-time certifier with slack $< \varepsilon_0/2$
would decide $\varepsilon_0$-Verify on the gap instances, contradicting
(3); a slack-$0$ bound decides the exact problem, requiring
$2^{\Omega(n)}$ under ETH. $\square$

\vspace{2pt}
\noindent\textbf{Remarks.}
\textit{(i) The single-layer case is polynomial.} For a \emph{single}
fused layer $\sigma(Wx + b)$, the supremum
$\sup_x |\sigma(Wx+b) - \tilde{\sigma}(Wx+b)|$ is a piecewise-linear
function of $x$ whose supremum over a polytope is the max of finitely
many linear programs---polynomial. The hardness in part~(2) lives in the
\emph{nested} composition of coupled gates (the product structure),
not in any single fused layer; earlier reductions to ``single-layer
NP-hardness'' were incorrect and are superseded here.
\textit{(ii) Earlier framing corrected.} The previous version compared the
tight bound (NP-hard for both classes) with an upper bound (polynomial
for both classes, since CROWN/interval bounds are polynomial-time for
MLPs~\cite{alphabetacrown2021,katz2019marabou}), which does not separate
the classes. The valid separation is in the $\varepsilon$-dimension as
stated above: the \emph{decision problem} is poly for SVNN and NP-hard
for coupled architectures, and SVNN uniquely attains
(polynomial time, zero slack) simultaneously.
\textit{(iii) Empirical demonstration (E65).} For a trained SVNN the
closed-form bound computes instantly at any $N$; brute-force supremum
enumeration over the input box grows as $3^n$ ($n$ input dimensions):
124$\times$ at $n=4$, 641$\times$ at $n=6$, 5{,}719$\times$ at $n=8$---the
separation is observed, not just proven.


\subsection{Compile-Aware Generalization Bound}
\label{sec:svnn-generalization}

The SVNN framework has so far addressed the \textit{compilation} side:
whether a trained model can be deployed with correctness guarantees. We now
establish the \emph{compile-aware} generalization bound: the object of the
guarantee is the \emph{deployed} (LUT-compiled) network $\bar{N}$, not the
trained network $N$. This unifies learning theory with LUT discretization
in a single two-scale bound, and yields a novel, directly actionable
\emph{resolution-matching law} for embedded deployment: the LUT memory
should scale as the $2k$-th root of the sample count.

\setcounter{theorem}{5}
\begin{theorem}[Compile-Aware Generalization]
\label{thm:generalization}
Let $N$ be an SVNN with $L$ affine layers and $L-1$ activation layers,
per-activation-layer Lipschitz product $\leq \gamma$ (Condition~3), trained
on $n$ i.i.d.\ samples; let $\bar{N} = \mathrm{Compile}(N, h)$ be the
deployed order-$k$ LUT version with cell width $h$. Let $\ell$ be the loss,
Lipschitz $L_\ell$, with outputs and targets bounded in $[-C, C]$, inputs in
a compact box $\Omega \subset [-B, B]^{d_{\text{in}}}$. Then with
probability at least $1-\delta$ over the sample:
\begin{equation}
R(\bar{N}) \;\leq\; \hat{R}(N) +
C_1 L_\ell \gamma^{L-1} \sqrt{L\, d_{\max}}\,
\sqrt{\frac{\log(2d_{\max}) + \log(1/\delta)}{n}}
+ C_2 \gamma^{L-1} L \bigl( c_k M_k h^k + \varepsilon_{\mathrm{fp}} \|\mathrm{scale}\| \bigr),
\label{eq:compile_aware}
\end{equation}
where $c_k = Q_k/k!$ is the sharp LUT constant
(Theorem~\ref{thm:sharp-lut}), $M_k = \max_f \|f^{(k)}\|_\infty$ over
activations, and $C_1, C_2$ are absolute constants. Equivalently, the
deployed risk decomposes as
\begin{equation}
R(\bar{N}) = \hat{R}(N) + \underbrace{O\!\left(\frac{\gamma^{L-1}}{\sqrt{n}}\right)}_{\text{gap}}
+ \underbrace{O\!\left(\gamma^{L-1} c_k M_k h^k\right)}_{\text{LUT bias}},
\label{eq:two_scale}
\end{equation}
with two-scale convergence: gap $\to 0$ at rate $n^{-1/2}$, bias $\to 0$
at rate $N^{-k}$ ($N = (b-a)/h$ LUT points).
\end{theorem}

\vspace{2pt}
\noindent\textit{Proof sketch.}
\textbf{(1) Bias term.} By Theorem~\ref{thm:sharp-lut}, each activation
layer $\ell$ contributes pointwise error $|e_\ell| \leq c_k M_k^{(\ell)}
h_\ell^k$; propagation through the remaining $L-\ell$ linear maps amplifies
by $\prod_{j>\ell} \|W_j\|_\infty \leq \gamma^{L-1}$ (Condition~3 gives the
per-activation-layer product $\leq \gamma^{L-1}$ for the remaining
activation layers; the Box-Continuation Lemma below validates that the
reachable intermediate signals lie within the LUT domains, so the $M_k$
constants apply without reachability computation). Floating-point
roundoff contributes $\varepsilon_{\mathrm{fp}}\|\mathrm{scale}\|$ per
layer. \textbf{(2) Gap term.} Standard Rademacher complexity of the SVNN
class with spectrally bounded weights
(Bartlett--Foster--Telgarsky-type chain), applying Talagrand's contraction
\emph{once per activation layer} ($L-1$ contractions)---this is the source
of the constant $C_1$ with no unexplained factor (the earlier ``factor~4''
arose from applying the contraction twice on one activation layer, and the
log-term arithmetic was in error: $3\sqrt{\log(2/\delta)/(2n)} = 0.0389$,
not $0.0117$, at $n = 10{,}971$, $\delta = 0.05$; and $2\rho\gamma^L B /
\sqrt{n} = 0.0038$, not $0.0019$). \textbf{(3) Decomposition} by the
triangle inequality: $R(\bar{N}) \leq R(N) + L_\ell \|N - \bar{N}\|_\infty$,
then $R(N) \leq \hat{R}(N) + \mathrm{gap}$. $\square$

\vspace{4pt}
\begin{lemma}[Box Continuation: LUT Domains Cover the Reachable Set]
\label{lem:box_continuation}
For an SVNN with order-$k$ LUTs, the LUT-compiled intermediate boxes
propagate \emph{exactly} through the piecewise-affine layers in
$O(L \cdot d_{\max}^2 \cdot N)$ time (scan the $O(N)$ cells intersecting
each box), and the true signals lie within the LUT boxes inflated by the
accumulated propagated errors. Hence the $M_k$ constants are validated on
the \emph{reachable set} with no reachability computation. For
compact-support B-spline KAN the activation domain is all of
$\mathbb{R}$, so the lemma is trivial; for activations with
domain-dependent curvature it is what makes the bias term of
Theorem~\ref{thm:generalization} sound. (Empirically: without this
condition, intermediate features can leave the LUT grid---the
E53 measurement of in-domain worst-case error $\approx 17$ at every
LUT density is exactly this phenomenon: the LUT's clamped-endpoint
behavior diverges from the exact B-spline's zero-extension outside the
grid, and the discrepancy is $N$-independent.)
\end{lemma}

\vspace{4pt}
\begin{corollary}[Resolution Matching: Optimal Compilation]
\label{cor:resolution_matching}
Minimizing the two-scale bound (Eq.~\ref{eq:two_scale}) over $h$
(equivalently $N$) balances the two terms:
\begin{equation}
h^{*} = \Theta\!\left( \Bigl(
\frac{\sqrt{L\, d_{\max}/n}}{c_k M_k L}\,
\frac{C_1 L_\ell}{C_2}
\Bigr)^{1/k} \right),
\qquad
N^{*} = \Theta\!\left( n^{1/(2k)} \cdot \bigl(c_k M_k L^2 d_{\max}^{-1/2}\bigr)^{1/k} \right).
\label{eq:resolution_matching}
\end{equation}
The qualitative law: \emph{deployed LUT memory should scale as the $2k$-th
root of the sample count} (fourth root for $k=2$). Overshooting wastes PLC
memory with zero statistical gain; undershooting destroys the learned
function. This is a compiler rule that no existing PAC bound states, and
it makes the compilation resolution a function of the dataset size at
compile time.
\end{corollary}

\vspace{4pt}
\begin{corollary}[Fixed-Depth Tradeoff]
\label{cor:depth_threshold}
For fixed $(n, h)$, the bound (Eq.~\ref{eq:compile_aware}) is
\emph{non-monotone} in $L$ when $\gamma < 1$: the gap grows by a factor
$\gamma \sqrt{(L+1)/L}$ and the bias by $\gamma (L+1)/L$ per added layer.
Depth strictly helps only if the training-risk reduction beats the gap
growth: $\Delta \hat{R} \geq (\gamma - 1)\, C\, \sqrt{L\, d/\gamma^{L-1}
\cdot n}$, equivalently $n \geq n^{*}(\gamma, L, \Delta \hat{R}) =
\Theta(\gamma^{2(L-1)} L d / (\Delta \hat{R})^2)$. The earlier claim that
``deeper SVNN networks generalize strictly better'' (Corollary~1 of the
previous version) is false: it ignored that the gap term amplifies with
depth while the log-term is depth-independent; the corrected statement is
a quantitative threshold, which is both true and more useful.
\end{corollary}

\vspace{4pt}
\noindent\textbf{Quantitative instantiation for the trained KAN student.}
For the trained KAN $[28,16,4]$ with $L=2$, $\gamma = 0.182$ (empirically
measured, {\S}\ref{sec:svnn-kan}), input domain $[-3,3]^{28}$ ($B=3$),
cross-entropy loss ($L_\ell = 1$), and $n=10{,}971$ training samples, the
corrected bound is
\begin{equation}
R(\bar{N}) \leq \hat{R}(N) + \frac{2 \cdot 1 \cdot 0.182 \cdot 3}{\sqrt{10{,}971}}
+ 3\sqrt{\frac{\log(2/0.05)}{2 \cdot 10{,}971}}
+ C_2 \cdot 0.182 \cdot 2 \cdot \tfrac{1}{8} M_2 h^2,
\label{eq:quant_gen_bound}
\end{equation}
i.e., gap $\leq 0.0038 + 0.0389 = 0.0427$ (corrected from the earlier
$0.0136$, whose arithmetic was inconsistent with its own formula) plus the
LUT bias $0.0455\, C_2 M_2 h^2$. The empirical generalization gap (training
vs.\ test accuracy) is $0.0\%$, consistent with the certificate. The
resolution-matching law (Corollary~\ref{cor:resolution_matching}) predicts
$N^{*} \asymp n^{1/4} \approx 10{,}971^{1/4} \approx 10.3$ for $k=2$,
matching the deployed $N = 15$ within $1.5\times$.

\vspace{4pt}
\noindent\textbf{Contrast with standard MLP.}
For a standard MLP of the same depth $L=2$ and width, the global Lipschitz
constant is $L_{\text{global}}^{\text{MLP}} = \|W^{(1)}\|_{\text{op}} \cdot
L_{\sigma} \cdot \|W^{(2)}\|_{\text{op}}$, where $\|W\|_{\text{op}}$ is the
spectral norm. There is no structural guarantee that this product is less than
1: in fact, weight-clipping or spectral normalization are required to enforce
contractivity, at the cost of reduced expressivity~\cite{golowich2018size}.
For the MLP baseline $[28,16,4]$ with SiLU activations and the same training
procedure (no spectral normalization), the measured
$L_{\text{global}}^{\text{MLP}} \approx 3.8$, giving a Rademacher complexity
term of $2 \cdot 3.8 \cdot 3 / \sqrt{10{,}971} \approx 0.217$---an order of
magnitude larger than the KAN's $0.0019$.

\vspace{4pt}
\begin{remark}[The Six SVNN Theorems]
\label{rem:six_theorems}
The six theorems and two propositions form a complete theoretical
architecture for the SVNN framework:
\begin{itemize}
    \item \textbf{Theorem~1} ({\S}\ref{sec:method}): \textit{Instantiation}---NeuroPLC
          compiles a specific KAN with a provable error bound $\varepsilon_{\text{DA}}$.
    \item \textbf{Theorem~2} ({\S}\ref{sec:svnn-conditions}):
          \textit{Sufficiency}---Conditions~1--2 guarantee design-time
          correctness for any SVNN architecture.
    \item \textbf{Theorem~3} ({\S}\ref{sec:method}):
          \textit{Greedy LUT allocation optimality}---the max-heap algorithm is
          minimax-optimal via exchange argument.
    \item \textbf{Theorem~4} ({\S}\ref{sec:method}):
          \textit{Probabilistic depth-scaling}---under the random-sign model,
          error grows only linearly with depth (martingale concentration).
    \item \textbf{Theorem~7} ({\S}\ref{sec:z3-verifiability}):
          \textit{Z3 NRA verifiability}---per-function Z3 queries are guaranteed
          UNSAT when $M_2 \cdot h^2 / 8 \leq C(3)$ (de~Boor guarantee).
    \item \textbf{Theorem~8} ({\S}\ref{sec:svnn-conditions}):
          \textit{SVNN Compositional Closure}---the SVNN class forms an algebraic
          monoid, enabling modular certification.
    \item \textbf{Theorem~9} ({\S}\ref{sec:da-optimality}):
          \textit{Sharp k-th order LUT bounds}---the de~Boor constant
          $M_2h^2/8$ generalizes to all orders with explicit sharp
          constants $c_k = Q_k/k!$ and an extremal family attaining them
          exactly; DA ($k=2$) is minimax-optimal among affine design-time
          schemes.
    \item \textbf{Theorem~10} ({\S}\ref{sec:wcet}):
          \textit{WCET Guarantee}---parameterized worst-case execution time bound
          of $22.67$~ms for KAN $[28,16,4]$ on S7-1200.
    \item \textbf{Theorem~A--E} ({\S}\ref{sec:characterization}--{\S}\ref{sec:noninterference}):
          Five king-level theorems---Compilable Frontier Characterization (A),
          DA Galois Connection (B), DA Tightness (C), IR Type Soundness (D),
          and Non-Interference (E)---completing the theoretical framework.
    \item \textbf{Theorem~5} ({\S}\ref{sec:svnn-necessity}):
          \textit{$\varepsilon$-separation}---the LUT-verification decision
          problem is polynomial for SVNN and NP-hard for coupled
          architectures (self-contained 3-SAT reduction); SVNN uniquely
          attains (polynomial time, zero slack).
    \item \textbf{Theorem~6} ({\S}\ref{sec:svnn-generalization}):
          \textit{Compile-aware learning theory}---the deployed risk
          decomposes as $\hat{R}(N) + O(\gamma^{L-1}/\sqrt{n}) +
          O(\gamma^{L-1}c_kM_kh^k)$, with the resolution-matching law
          $N^{*} \asymp n^{1/(2k)}$ (Corollary~\ref{cor:resolution_matching}).
\end{itemize}
\noindent\textbf{Propositions} (structural membership):
\begin{itemize}
    \item \textbf{Proposition~1} ({\S}\ref{sec:svnn-mlp}): Standard MLPs
          \textit{do not} satisfy the SVNN tight-bound structure.
    \item \textbf{Proposition~2} ({\S}\ref{sec:svnn-chebykan}): ChebyKAN
          \textit{does} satisfy SVNN Conditions~1--2, 496/512 Z3-verifiable.
    \item \textbf{Proposition~9} ({\S}\ref{sec:c2bv-family}): Operation
          Separation Principle---the $C^2$-BV family (B-spline, Fourier,
          Wavelet, RBF, Chebyshev) is universally SVNN-compliant.
    \item \textbf{Propositions~3--8}: DA Segment-Exactness, ChebyKAN SVNN,
          Fine-Tuning Stability, IR Minimality, Compiler Complexity,
          and Optimization Pass Soundness (see respective sections).
\end{itemize}
\vspace{4pt}
\noindent\textbf{King-Level Theorems A--E} (abstract interpretation + formal methods):
\begin{itemize}
    \item \textbf{Theorem~A} ({\S}\ref{sec:characterization}): Compilable
          Frontier Characterization---SVNN Conditions~1--2 are necessary
          and sufficient for dimension-optimal affine certification.
    \item \textbf{Theorem~B} ({\S}\ref{sec:galois}): DA Galois Connection---
          DA is a canonical instance of Cousot abstract interpretation.
    \item \textbf{Theorem~C} ({\S}\ref{sec:tightness}): DA Tightness---the
          $M_2 \cdot h^2 / 8$ bound is semantically attained.
    \item \textbf{Theorem~D} ({\S}\ref{sec:ir-semantics}): IR Type Soundness---
          compiler correctness as a type-theoretic result.
    \item \textbf{Theorem~E} ({\S}\ref{sec:noninterference}): SVNN
          Non-Interference---first formal proof of safe neural inference
          insertion in certified PLC programs.
\end{itemize}
Together, these theorems characterize the SVNN framework across four
dimensions: \textit{what} correctness guarantee is achievable (T1, T3, T4),
\textit{which architectures} admit it (T2, P1, P2), \textit{why} the
conditions cannot be weakened (T5), and \textit{what learning benefit}
follows from the strongest condition (T6).
\end{remark}
